\section{The main theorem}

\begin{theorem}[Alpöge's formal proof]\label{thm:grh}
Let \(q\ge1\) be an integer and let \(\chi\) be a Dirichlet character
modulo \(q\).
\begin{enumerate}
\item If \(\RePart s>1/2\), then \(L(s,\chi)\ne0\), with the exception
      of \(s=1\) when \(\chi\) is principal.
\item If \(0<\RePart s<1\) and \(L(s,\chi)=0\), then \(\RePart s=1/2\).
\item All zeros of the Riemann zeta function other than the trivial
      zeros \(-2,-4,\ldots\) lie on \(\RePart s=1/2\).
\end{enumerate}
\end{theorem}

\noindent
Here \(L(s,\chi)\) is the meromorphic continuation of
\(\sum_{n\ge1}\chi(n)n^{-s}\) from \(\RePart s>1\).
The character \(\chi\) may be imprimitive.
The \(L\)-function of the character modulo \(1\) is \(\zeta(s)\).
Of the five formal statements, two express (3).
The other three express (1), (2), and the specialization of (2)
to primitive characters.
See the \leansrc{Statements}{formal statements} and the
\leansrc{Main}{main declarations}.

We first deduce (1) from a one-sided bound for a weighted prime-power sum
(Section~\ref{sec:analytic}). We then place weights at the infinite places
of an auxiliary number field to express the remainder in terms of
divisibility of theta derivatives (Section~\ref{sec:data}).
Section~\ref{sec:lattices} combines the infinite and finite places for
a single Gram determinant. Sections~\ref{sec:upper-bound}--\ref{sec:determinant-estimates}
establish the comparison's upper bound, dimension estimate, and bound
for the determinant ratios under twisting. Section~\ref{sec:conclusion}
deduces the theorem.

The proof follows Alpöge's announcement \cite{alpoege} and the
\texttt{grh/} project, which the source repository calls its readable
edition \cite{grh-source}. References beginning with L point to the
corresponding declarations in the appendix. The proofs of the estimates
for theta valuations, the Poincaré second moment, and the choice of
twists and projection denominators appear in their respective sections.
For auxiliary results on the Weil representation and local orbit sums,
we give the forms used and cite the corresponding sources. Blue notes are editorial
explanations. The overview that follows describes the main ideas and
points to related accounts.
